% 第五章 系统详细设计与实现
% 本章在第四章总体设计基础上，重点说明核心功能的实现方法。
% 已按“减少普通业务算法、突出视频上传转码与 AI 审核实现”的方向进行压缩整理。

\chapter{系统详细设计与实现}

用户端功能模块是普通用户访问系统的主要入口，主要承担视频浏览、视频搜索、视频播放、用户认证和互动操作等功能。该模块并不是单一的数据读取模块，而是由主站业务服务、互动服务、缓存组件、检索组件和数据库共同协作完成。公开视频浏览、分类查询和视频搜索允许未登录用户访问，而点赞、收藏、投币、关注、评论和发送弹幕等操作需要用户完成登录认证后才能执行。

\section{用户端功能模块实现}

用户进入系统后，可以通过首页浏览公开视频内容，也可以根据分类或关键词搜索目标视频；当用户进入视频播放业务后，系统需要加载视频详情、分 P 文件、播放资源、弹幕数据以及互动状态；当用户执行评论、弹幕、点赞、收藏或投币等操作时，系统需要进一步校验登录状态，并将用户行为记录到对应的数据表中。

从服务协作角度看，用户端功能主要由主站业务服务和互动服务共同完成。主站业务服务负责账号认证、视频分类、视频列表、视频详情和播放文件等核心业务；互动服务负责点赞、收藏、投币、评论、弹幕和关注等用户行为；Redis 负责验证码、登录态和搜索热词等临时数据；Elasticsearch 负责视频关键词检索；MySQL 负责用户、视频、分类、评论、弹幕和互动行为等核心数据的持久化存储。

% 此处放图5-1：用户端功能模块协作图
% 建议图中包含：用户访问请求、网关服务、主站业务服务、互动服务、MySQL、Redis、Elasticsearch。
% 图名建议：用户端功能模块协作图

用户端功能模块的实现重点在于区分公开访问业务和登录后业务。公开访问业务主要包括首页视频浏览、分类筛选、关键词搜索和公开视频详情查询，这些功能不要求用户必须登录，系统只需要根据视频状态过滤出审核通过且已发布的视频数据。登录后业务主要包括点赞、收藏、投币、关注、评论和发送弹幕等操作，这些行为必须绑定具体用户身份，因此系统在执行业务前需要先校验登录态。

为了保证用户访问过程中的数据一致性，系统在各业务模块中采用统一的状态校验方式。例如，视频播放前需要判断视频是否存在、是否审核通过、是否已经发布以及播放文件是否已经生成；用户执行互动操作前需要判断用户是否登录、账号状态是否正常以及是否已经存在相同行为记录。通过这些校验逻辑，系统能够避免非法视频被访问、重复行为被写入以及未登录用户执行身份绑定操作等问题。

\subsection{用户端整体功能实现}

用户端整体功能实现按照“请求接入、身份判断、业务分发、数据处理、结果返回”的流程展开。当用户发起访问请求后，系统首先由网关服务接收请求，并根据请求路径转发到对应的业务服务。对于公开视频浏览、分类查询、视频搜索和视频详情查询等公开接口，系统可以直接进入业务处理流程；对于点赞、收藏、投币、关注、评论和弹幕发送等需要用户身份的接口，系统需要先读取登录令牌，并从 Redis 中查询登录态是否有效。

在业务分发过程中，系统根据不同请求类型调用不同的业务模块。账号认证请求由账号认证服务处理，首页浏览和搜索请求由视频业务服务处理，视频播放请求由视频详情和播放文件服务处理，评论、弹幕和互动行为请求由互动服务处理。各业务服务在完成参数校验后，分别访问 MySQL、Redis 或 Elasticsearch 获取数据，并将结果封装为统一响应对象返回。

用户端整体功能的数据流转可以概括为以下几个步骤。首先，系统接收用户访问请求，并判断该请求是否需要登录认证。其次，若请求需要登录，则根据请求中的登录标识读取 Redis 中的登录态；若登录态不存在或已过期，则直接返回未登录结果。再次，若身份校验通过，则业务服务根据请求参数查询或写入业务数据。最后，系统将处理结果封装后返回给用户访问入口。

在数据对象设计上，用户端整体功能主要涉及用户认证对象、视频摘要对象、视频详情对象、播放文件对象、评论对象、弹幕对象和用户行为对象。用户认证对象用于保存用户编号、昵称、头像和登录状态等信息；视频摘要对象用于首页和搜索结果展示；视频详情对象用于视频播放业务；播放文件对象用于描述视频分 P 文件和播放资源；评论对象和弹幕对象用于社区互动；用户行为对象用于记录点赞、收藏、投币和关注等行为。

整体来看，用户端功能实现并不是孤立的单接口处理，而是围绕视频内容消费过程形成完整业务闭环。首页浏览和搜索帮助用户发现视频，播放页完成内容消费，认证模块保证用户身份可信，互动模块记录用户行为。多个模块共同协作，使系统能够同时支持游客访问和登录用户互动。

\subsection{首页浏览与视频搜索实现}

首页浏览与视频搜索功能主要负责视频内容发现。用户进入系统后，可以直接浏览公开视频列表，也可以按照分类筛选视频，或者通过关键词搜索目标视频。该功能的核心是根据不同访问条件构造查询请求，并从公开视频数据中筛选出符合条件的视频结果。

首页浏览功能主要处理分类查询和公开视频列表查询。系统首先读取视频分类信息，分类数据用于组织一级分类和二级分类。用户选择分类后，系统根据分类编号构造查询条件，从视频信息表中查询已经审核通过并正式发布的视频记录。查询结果通常需要包含视频编号、视频标题、封面、作者信息、播放量、弹幕数量和发布时间等摘要信息，以便用户快速判断视频内容。

视频搜索功能在普通视频列表查询的基础上增加了关键词检索。用户提交关键词后，系统会先对关键词进行合法性校验，避免空关键词或无效关键词进入检索流程。校验通过后，系统将关键词提交给检索服务，由 Elasticsearch 根据视频标题、简介和标签等索引字段执行全文检索。同时，系统可以将关键词写入 Redis，用于统计热门搜索词，为后续首页热词展示提供数据来源。

% 此处放图5-2：首页浏览与视频搜索流程图
% 建议图中包含：用户请求、分类浏览、关键词搜索、MySQL视频查询、Elasticsearch检索、Redis热词统计、视频摘要结果返回。
% 图名建议：首页浏览与视频搜索流程图

首页浏览和视频搜索虽然都返回视频列表，但二者的数据来源和查询方式存在差异。分类浏览主要依赖 MySQL 中的视频分类字段和视频状态字段，适合按照分类、发布时间或推荐状态查询公开视频；关键词搜索主要依赖 Elasticsearch 的全文检索能力，适合根据用户输入内容快速匹配视频标题、简介和标签。两种方式最终都会将结果封装为视频摘要数据返回，从而保证返回结构统一。

在实现过程中，系统需要特别注意公开视频状态控制。由于视频从投稿到发布需要经过上传、转码、审核和发布等多个阶段，因此首页和搜索结果中不能直接展示投稿中的视频数据。系统必须过滤未审核、审核未通过、转码失败或状态异常的视频，只返回已经正式发布且资源可访问的视频。这样可以保证用户在首页和搜索结果中看到的视频均处于可播放状态。

首页浏览与视频搜索模块的数据处理过程主要包括查询条件构造、数据源选择、结果过滤和结果封装四个步骤。查询条件构造用于生成分类、分页、排序或关键词参数；数据源选择用于判断当前请求应访问 MySQL 还是 Elasticsearch；结果过滤用于排除状态异常的视频；结果封装用于将数据库实体或检索结果转换为视频摘要对象。通过该流程，系统能够同时支持分类浏览和关键词搜索两种内容发现方式。

\subsection{视频播放页实现}

视频播放页功能是用户端内容消费的核心。用户从首页、搜索结果或相关视频入口进入某个视频后，系统需要根据视频编号加载视频详情、分 P 文件、播放资源、作者信息、弹幕数据和互动状态。该功能不是简单地返回视频文件地址，而是需要围绕视频状态、播放资源、用户行为和社区互动进行综合处理。

视频播放业务首先需要校验视频状态。系统根据视频编号查询视频主数据，判断视频是否存在、是否已经审核通过、是否处于正常发布状态。如果视频不存在、未发布或状态异常，则系统不允许继续加载播放资源。只有当视频状态合法时，系统才会继续查询对应的分 P 文件和播放资源。

在分 P 文件处理方面，一个视频可以包含多个播放文件，每个文件对应一个分 P。系统需要按照文件序号查询并返回当前视频的分 P 文件列表。默认情况下，系统可以选择第一个可播放文件作为初始播放文件；当用户切换分 P 时，系统根据文件编号重新读取对应播放资源，并加载该分 P 下的弹幕数据。由于弹幕与具体播放文件和播放时间点绑定，因此分 P 切换时必须重新加载弹幕，避免不同分 P 的弹幕数据混淆。

% 此处放图5-3：视频播放状态流转图
% 建议图中包含：请求视频详情、校验视频状态、查询分P文件、加载播放资源、加载弹幕数据、正常播放、切换分P、提交互动行为、视频不可播放、无播放资源。
% 图名建议：视频播放状态流转图

视频播放业务还需要处理用户互动状态。如果用户未登录，系统只返回视频基础信息、播放文件和公开统计数据；如果用户已经登录，系统还需要查询当前用户是否已经点赞、收藏、投币以及是否关注作者。通过返回这些个性化状态，系统可以避免重复互动操作，并为后续用户行为处理提供依据。

弹幕功能是视频播放页中的重要互动功能。系统加载弹幕时，需要根据视频编号和文件编号查询弹幕数据，并按照播放时间排序。用户发送弹幕时，系统需要校验登录状态、视频状态、文件编号和弹幕内容，校验通过后将弹幕文本、用户编号、视频编号、文件编号、播放时间、颜色和显示模式等数据写入弹幕表。弹幕保存成功后，系统可更新视频弹幕统计数量。

评论功能与弹幕功能不同，评论围绕整个视频展开，不绑定具体播放时间点。系统查询评论时，主要根据视频编号读取评论列表，并根据父评论编号组织回复关系。用户发表评论或回复时，同样需要登录校验和内容校验。评论保存后，系统可以更新视频评论统计或评论回复数量。

点赞、收藏和投币等互动行为通过统一的用户行为处理逻辑实现。系统根据用户编号、视频编号和行为类型判断当前行为是否已经存在，避免重复点赞、重复收藏或重复投币。对于投币行为，还需要进一步校验用户可用资源是否满足操作要求。行为处理成功后，系统同时更新用户行为记录和视频统计数据，使播放页中的互动数量能够反映最新状态。

综上，视频播放页实现的重点在于保证视频内容状态合法、播放资源可用、分 P 文件清晰、弹幕数据准确以及用户互动行为可追踪。该模块将视频主数据、播放文件、弹幕评论和用户行为整合在同一业务链路中，是用户端功能模块中业务关联度最高的部分。

\subsection{用户认证与登录状态实现}

用户认证与登录状态实现是用户端业务的基础支撑。系统允许游客访问公开视频内容，但对于评论、弹幕、点赞、收藏、投币、关注以及用户中心等身份绑定业务，必须要求用户处于登录状态。因此，认证模块需要完成用户身份校验、登录态保存、自动登录检查和退出登录处理。

用户登录流程首先从验证码生成开始。系统生成验证码文本和验证码标识后，将验证码结果保存到 Redis，并设置有效期。用户提交登录请求时，需要同时提交账号、密码、验证码和验证码标识。系统先根据验证码标识从 Redis 中读取验证码结果并进行比对，验证码校验通过后，再查询用户信息并校验密码和账号状态。只有验证码、账号密码和账号状态均校验通过后，系统才会生成登录令牌。

登录成功后，系统将用户编号、昵称、头像、账号状态和登录令牌等信息保存到 Redis 中，并设置登录过期时间。后续用户访问需要登录的接口时，系统不需要重复校验账号密码，而是根据请求中的登录标识到 Redis 中读取登录态。如果 Redis 中存在有效登录信息，则说明用户处于登录状态；如果登录信息不存在或已过期，则系统返回未登录或登录失效结果。

% 此处放图5-4：登录状态流转图
% 建议图中包含：未登录、获取验证码、提交登录信息、登录失败、已登录、自动登录检查、执行登录后操作、退出登录、登录过期。
% 图名建议：登录状态流转图

自动登录用于解决用户刷新页面或重新访问系统时登录状态恢复的问题。当用户再次访问系统时，系统会读取请求中携带的登录标识，并到 Redis 中查询对应登录态。如果登录态有效，则重新返回用户认证信息；如果登录态不存在或已经过期，则保持未登录状态。通过自动登录机制，系统能够减少用户重复登录次数，同时保证登录状态仍然受 Redis 有效期控制。

退出登录流程则用于主动清除用户登录状态。用户发起退出请求后，系统根据登录令牌删除 Redis 中保存的登录信息，使该令牌立即失效。退出成功后，用户再次访问需要登录的接口时，即使请求中仍携带旧令牌，也无法通过登录态校验，从而保证退出行为生效。

认证模块还需要承担业务访问控制职责。对于公开视频浏览、分类查询、搜索和视频详情读取等公开业务，系统不强制登录；对于点赞、收藏、投币、关注、评论和弹幕发送等身份绑定业务，系统必须在业务执行前完成登录校验。通过统一的认证校验逻辑，可以避免不同业务模块重复编写登录判断代码，也能减少权限控制不一致的问题。

用户认证与登录状态实现采用“验证码校验、账号密码校验、Redis 保存登录态、接口访问前校验身份”的方式。该方式既支持游客访问公开视频内容，又能保证用户互动行为与具体用户身份绑定，为用户中心、创作中心和互动功能提供了统一可靠的身份基础。

\section{创作中心与视频投稿实现}

创作中心与视频投稿模块面向内容创作者提供稿件管理、视频资源上传、投稿信息保存和稿件状态维护等功能。与普通视频浏览功能不同，视频投稿涉及用户身份校验、文件分片上传、投稿主数据保存、分 P 文件绑定和转码任务初始化等多个环节。为了保证投稿流程清晰，系统将视频投稿过程划分为资源上传阶段和投稿保存阶段：资源上传阶段只负责将视频分片保存到临时目录并维护上传状态；投稿保存阶段才真正生成视频投稿主记录和分 P 文件记录，并将需要处理的文件加入转码队列。

创作中心的业务处理以当前登录用户为基础。创作者进行稿件查询、投稿保存、投稿详情读取、互动设置修改和稿件删除等操作时，系统都需要先从登录态中解析用户编号，再根据用户编号限制数据访问范围。这样可以保证创作者只能管理自己的稿件，避免不同用户之间的视频投稿数据被越权访问或修改。

从数据结构看，视频投稿过程主要涉及用户对象、投稿视频主对象、投稿视频文件对象、上传任务对象和转码任务对象。用户对象表示投稿行为的发起者；投稿视频主对象保存视频标题、封面、分类、标签、简介、投稿类型、互动设置和投稿状态；投稿视频文件对象保存每个分 P 文件的上传标识、文件名称、文件序号、文件路径和转码结果；上传任务对象保存分片上传过程中的临时状态；转码任务对象则由投稿文件记录生成，用于后续资源服务处理。

% 此处放图5-X：投稿核心对象关系图
% 建议图中包含：UserInfo、VideoInfoPost、VideoInfoFilePost、UploadTask、TranscodeTask。
% 建议关系：
% UserInfo 1 对多 VideoInfoPost；
% VideoInfoPost 1 对多 VideoInfoFilePost；
% VideoInfoFilePost 关联 UploadTask；
% VideoInfoFilePost 生成 TranscodeTask。
% 图名建议：投稿核心对象关系图

该模块的设计重点在于区分投稿态数据和正式发布数据。投稿态数据用于保存创作者正在编辑、转码、待审核或审核失败的视频；正式发布数据只保存审核通过后可公开访问的视频。通过这种数据分层方式，系统可以避免未审核或处理失败的视频进入首页列表、搜索结果和公开视频播放流程，同时也便于创作者在审核前后对稿件进行独立管理。

\subsection{创作中心功能结构}

创作中心功能结构围绕创作者稿件生命周期展开，主要包括稿件列表查询、稿件状态统计、投稿详情读取、投稿信息保存、互动设置维护和稿件删除等功能。创作者进入创作中心后，可以查看自己提交过的稿件，并按照状态筛选正在处理、审核通过或审核失败的视频；在新增或编辑投稿时，系统保存投稿主数据和分 P 文件信息；当稿件已经发布后，创作者可以根据业务规则修改部分互动设置或删除对应稿件。

稿件列表查询以当前登录用户为约束条件。系统从登录态中解析用户编号后，只查询该创作者自己的投稿记录。查询时可以根据稿件状态、页码和视频名称等条件进行筛选，并按照创建时间或更新时间返回结果。该处理方式既满足创作者管理稿件的需要，也保证了投稿数据的访问范围可控。

稿件状态统计用于展示创作者不同状态下的视频数量。系统按照投稿状态对当前用户的稿件进行统计，例如处理中、审核通过和审核失败等状态。其中，处理中状态通常表示视频尚未进入最终发布结果，可能处于转码中、待审核或其他中间状态。通过状态统计，创作者可以快速了解当前投稿处理情况，判断哪些视频仍处于系统处理阶段，哪些视频已经发布，哪些视频需要根据失败原因重新修改。

投稿详情读取用于支持稿件编辑。系统根据视频编号查询投稿主表，并校验该视频是否属于当前登录用户；如果视频不存在或不属于当前用户，则返回资源不存在或无权限结果。校验通过后，系统继续查询该视频对应的分 P 文件记录，并按照文件序号升序排列，最终将视频主信息和分 P 文件列表一起返回。该过程保证创作者在编辑稿件时能够恢复完整的投稿内容和文件结构。

互动设置维护用于控制视频发布后的互动能力，例如是否允许评论、是否允许弹幕等。该类设置属于视频业务规则的一部分，需要与视频主数据绑定保存。系统在保存互动设置时，同样需要校验视频归属关系，保证只有视频创作者本人才能修改对应设置。

稿件删除功能需要结合稿件状态进行处理。系统需要先判断当前登录用户是否为稿件所有者，再根据稿件状态决定是否允许删除。对于正在转码或审核中的稿件，直接删除可能影响异步任务执行，因此需要结合业务状态进行限制或延迟资源清理。对于可以删除的稿件，系统不仅需要删除投稿主数据和分 P 文件数据，还需要处理对应的资源文件清理任务。

整体来看，创作中心功能结构以“稿件”为核心对象展开，围绕稿件的创建、编辑、查询、统计和删除形成完整管理链路。该模块不负责视频转码的具体执行，也不负责公开视频播放消费，而是负责将创作者提交的视频信息和上传资源组织成可进入后续处理流程的投稿数据。

\subsection{视频分片上传实现}

视频文件通常体积较大，如果采用普通文件表单一次性提交，容易受到请求超时、网络中断和服务端单次请求大小限制的影响。因此，系统在资源服务中采用分片上传方式处理视频文件。分片上传将一个完整视频拆分为多个小块依次提交，服务端在每次收到分片后保存到临时目录，并在 Redis 中更新上传进度和文件大小等状态信息。

视频分片上传流程首先由上传前初始化开始。创作者选择视频文件后，系统向资源服务提交文件名称和分片数量。资源服务根据当前登录用户、原始文件名和分片数量生成上传标识，并将上传临时信息保存到 Redis 中。上传标识是后续分片上传和投稿保存之间的重要关联字段，用于标识某个用户的一次视频上传任务。

% 此处放图5-X：视频分片上传流程图
% 建议图中包含：申请上传标识、生成 uploadId、Redis 保存上传任务、循环上传分片、校验 uploadId、校验 chunkIndex、写入临时目录、更新 Redis 上传进度、上传完成、投稿保存时引用 uploadId。
% 图名建议：视频分片上传流程图

分片上传时，每个上传请求需要携带上传标识、当前分片索引和分片文件。资源服务首先从登录态中读取用户编号，再根据用户编号和上传标识从 Redis 中获取上传临时信息。如果 Redis 中不存在对应记录，说明上传标识无效或上传任务已经过期，系统会拒绝本次上传。该校验可以防止用户伪造上传标识，也可以避免服务端接收无归属的临时文件。

在接收分片之前，系统还需要进行文件大小和分片顺序校验。文件大小校验用于防止上传文件超过系统配置中的视频大小限制；分片顺序校验用于判断当前分片索引是否合法，避免出现跳片上传或越界上传。校验通过后，系统将当前分片写入该上传任务对应的临时目录，并以分片索引作为临时文件名保存。随后系统更新 Redis 中的当前分片索引和累计文件大小，使服务端能够记录上传进度。

如果创作者取消上传或需要重新选择文件，系统可以根据上传标识删除上传临时信息，并清理对应临时目录。删除上传任务时，系统同样需要校验登录用户和上传标识的对应关系，只有当前用户自己的上传任务才能被删除。这样可以避免用户误删其他用户的临时资源。

需要注意的是，分片上传阶段只负责将文件分片保存到临时目录，并不会立即合并文件，也不会生成正式播放资源。视频能否进入后续转码和审核流程，取决于投稿保存阶段是否将该上传标识写入分 P 文件列表。也就是说，上传成功的文件只是投稿资源的临时准备结果，只有当创作者提交完整投稿信息后，系统才会将分 P 文件与视频投稿主记录绑定。

\subsection{投稿信息保存与状态初始化实现}

投稿信息保存是视频投稿流程中的核心步骤。创作者完成封面上传、视频分片上传和基础信息填写后，系统将视频标题、封面、一级分类、二级分类、投稿类型、标签、简介、互动设置以及分 P 文件列表一起提交到主站业务服务。主站业务服务在接收到请求后，首先解析登录用户信息，然后将提交的分 P 文件列表转换为投稿文件对象，最后调用投稿业务服务完成保存。

投稿保存过程首先进行前置校验。系统会根据配置检查分 P 数量是否超过允许上限；如果本次请求携带视频编号，则说明该操作属于编辑已有稿件，系统需要先查询对应投稿记录是否存在，并判断该视频当前状态是否允许修改。对于正在转码或正在审核的视频，系统不允许再次编辑，避免异步任务处理过程中数据被修改，导致转码结果、审核结果和投稿信息不一致。

% 此处放图5-X：投稿信息保存与状态初始化流程图
% 建议图中包含：提交投稿信息、校验登录用户、校验分P数量、判断新增或编辑、保存 video_info_post、保存 video_info_file_post、初始化视频状态、初始化文件转码状态、加入 Redis 转码队列。
% 注意：该图只画到“加入转码队列”为止，不要继续画资源服务合并与转码，避免和下一节重复。
% 图名建议：投稿信息保存与状态初始化流程图

对于新增投稿，系统会生成新的视频编号，并设置创建时间、更新时间和初始状态。新增视频的初始状态通常为转码中，表示该视频已经完成投稿信息保存，但其上传文件尚未完成合并、转码和播放资源生成。随后系统将视频主信息写入 \texttt{video\_info\_post} 表，并为每个分 P 文件生成文件编号、设置文件序号、绑定视频编号和用户编号，同时将文件转码结果设置为转码中状态。

对于编辑投稿，系统需要先查询数据库中已有的分 P 文件记录，并与本次提交的分 P 文件列表进行对比。若数据库中存在而本次提交列表中不存在的文件，说明该分 P 已被创作者删除，系统需要删除对应的投稿文件记录，并将其文件路径加入待删除队列。若本次提交的文件没有文件编号，说明该文件属于新增分 P，需要生成新的文件编号并加入转码队列。若只是标题、简介、标签或封面等基础信息发生变化，则不需要重新转码，但需要重新进入审核流程。

投稿主信息的状态初始化与文件变化情况有关。如果是新增视频，或者编辑时新增了分 P 文件，则系统将视频状态设置为转码中，因为新增文件需要经过资源处理后才能播放。如果没有新增文件，但视频标题、封面、标签、简介或投稿类型等基础信息发生变化，则系统将视频状态设置为待审核，表示资源已经存在，但内容信息需要重新审核。

投稿文件保存后，系统将新增的分 P 文件加入 Redis 转码队列。到此为止，创作中心与视频投稿模块完成了投稿业务数据的保存和状态初始化工作。后续的文件合并、视频转码、HLS 播放资源生成和转码结果回写由资源服务异步完成，不再属于本节的主要内容。

投稿信息保存与状态初始化的本质，是将“上传完成的临时资源”和“可处理的视频稿件”建立关联。分片上传阶段只保存文件临时状态，而投稿保存阶段才真正生成视频投稿主记录和分 P 文件记录。系统通过 \texttt{uploadId} 将上传任务与分 P 文件绑定，通过 \texttt{videoId} 将投稿主信息与多个分 P 文件绑定，通过投稿状态字段描述视频当前所处的业务阶段。

该设计可以有效支持新增投稿和编辑投稿两类场景。新增投稿时，系统从无到有创建视频主记录、文件记录并初始化转码状态；编辑投稿时，系统根据差异判断是否需要重新转码或重新审核。通过投稿态表、文件态表、上传标识和转码队列的配合，系统能够完整描述视频从提交到进入资源处理队列之前的业务过程，为后续资源服务处理和审核流程提供可靠的数据基础。

\section{视频转码与资源服务实现}

视频转码与资源服务模块负责处理投稿完成后的视频资源加工、播放资源生成和资源访问等功能。上一节已经说明了视频分片上传、投稿保存和转码任务入队过程，因此本节不再重复描述上传阶段和投稿阶段，而是重点说明资源服务如何从 Redis 转码队列中消费任务，并将已经提交的视频文件处理成可访问的 HLS 播放资源。

视频文件在投稿保存后并不会立即变成可播放资源。系统只是将新增分 P 文件加入转码队列，并将视频和文件状态初始化为转码中。资源服务中的后台任务会持续监听 Redis 转码队列，当队列中存在待处理文件时，资源服务取出对应的投稿文件对象，根据用户编号、上传标识、视频编号和文件编号执行临时文件迁移、分片合并、视频时长读取、编码兼容处理和 HLS 资源生成。处理完成后，资源服务再将结果回写给主站业务服务，由主站服务更新分 P 文件状态和视频投稿状态。

该模块的核心设计思想是将耗时的资源处理过程与投稿接口解耦。投稿接口只负责保存业务数据和初始化任务，不等待视频合并与转码完成；资源服务通过异步任务消费队列并处理文件，处理结果通过回写接口反馈给主站服务。这样可以避免用户提交投稿时长时间阻塞，也能够提升系统对大文件和多分 P 投稿的处理能力。

\subsection{资源服务职责}

资源服务主要承担文件资源处理和播放资源输出职责。从功能边界看，资源服务不负责视频标题、分类、标签、简介等投稿业务信息维护，也不直接负责视频审核逻辑，而是围绕文件资源本身提供加工和读取能力。其职责主要包括转码任务消费、上传临时文件迁移、视频分片合并、视频编码处理、HLS 播放资源生成、播放资源访问和播放信息记录等。

首先，资源服务负责消费转码队列。投稿信息保存成功后，主站服务将新增分 P 文件加入 Redis 转码队列。资源服务启动后，后台任务线程会持续从队列中读取待处理的投稿文件对象。若队列为空，任务线程短暂休眠后继续轮询；若读取到任务，则调用资源处理组件执行文件处理逻辑。该方式使资源服务能够持续处理投稿产生的视频文件任务，并且不会阻塞用户请求。

其次，资源服务负责将上传阶段产生的临时文件转换为正式资源。资源服务根据任务中的用户编号和上传标识，从 Redis 中读取上传任务信息，并定位上传临时目录。随后系统将临时目录迁移到正式视频资源目录，并删除 Redis 中的上传临时信息。完成迁移后，资源服务按照分片索引合并文件，生成完整视频文件。

再次，资源服务负责视频转码和播放资源生成。完整视频生成后，系统调用 FFmpeg 工具读取视频时长和编码格式；如果视频编码格式存在兼容性问题，则先进行编码转换；随后生成 HLS 播放所需的 \texttt{m3u8} 索引文件和 \texttt{ts} 切片文件。该处理过程使大体积视频能够以流式方式被播放器加载，提升播放稳定性和访问体验。

最后，资源服务负责播放资源访问。用户播放视频时，并不直接访问服务器真实文件路径，而是根据分 P 文件编号向资源服务请求播放索引和切片文件。资源服务根据文件编号查询正式文件记录或投稿态文件记录，定位资源目录后读取对应文件并写入响应。对于正式发布视频，资源服务还会记录播放信息，用于后续异步统计播放量。

资源服务与主站服务之间采用明确的职责分离。主站服务负责保存视频业务状态和投稿状态，资源服务负责处理真实文件资源；资源服务完成转码后，不直接操作主站数据库，而是通过服务调用回写转码结果。这样可以避免不同服务直接共享数据库操作逻辑，也使视频资源处理过程具有较好的独立性。

\subsection{视频文件合并与转码实现}

视频文件合并与转码是资源服务的核心处理过程。系统在分片上传阶段将视频拆分为多个分片文件保存到临时目录，而后续播放和审核需要完整的视频资源。因此，在进入转码前，资源服务首先需要根据上传标识找到对应的临时目录，并将临时目录中的分片文件迁移到正式视频资源目录。

转码任务由资源服务中的后台任务触发。系统启动后，任务线程持续从 Redis 转码队列中获取待处理文件对象。如果队列为空，线程短暂休眠后继续轮询；如果获取到待处理文件，则调用转码组件处理该文件。通过这种方式，系统可以持续处理投稿产生的视频文件任务，并且不会阻塞用户投稿请求。

% 此处放图5-X：视频文件合并与转码流程图
% 建议图中包含：从 Redis 转码队列取任务、读取上传临时信息、复制临时目录到正式目录、删除临时目录、合并分片、读取视频时长、判断编码格式、HEVC 转 H264、生成 HLS 切片、回写转码结果。
% 图名建议：视频文件合并与转码流程图

在文件合并过程中，资源服务首先从 Redis 中读取上传任务信息，根据上传任务中的文件路径定位临时分片目录。随后系统将临时目录复制到正式视频目录，并删除临时目录和 Redis 中的上传临时信息，表示该上传任务已经进入正式处理阶段。迁移完成后，系统按照分片索引依次读取分片文件，并通过随机访问文件方式将多个分片追加写入到一个完整视频文件中。合并完成后，如果配置要求清理源分片，系统会删除已经合并的分片文件，只保留完整视频和后续生成的播放资源。

完整视频生成后，资源服务需要读取视频元信息。系统调用 FFmpeg 工具中的 \texttt{ffprobe} 命令获取视频时长，并将时长和完整文件大小写入待回写对象。视频时长用于播放页展示、审核信息展示和视频总时长统计；文件大小用于记录资源占用情况，也可以为后续资源管理和清理提供依据。

在编码处理方面，系统会先读取视频编码格式。如果视频编码为 HEVC，由于部分播放环境对 HEVC 的兼容性较差，资源服务会先将其转换为 H.264 编码的视频文件，再继续执行 HLS 切片处理。对于普通 H.264 视频，则可以直接进入切片阶段。通过编码判断和兼容转换，可以降低播放端无法正常解码的风险。

HLS 切片生成是转码处理的最后阶段。系统先将完整视频转换为中间 \texttt{ts} 文件，再按照固定时间间隔对视频进行切片，并生成对应的 \texttt{m3u8} 索引文件。切片完成后，播放目录中会包含一个索引文件和多个 \texttt{ts} 切片文件。播放器访问时先读取索引文件，再根据索引文件中的切片信息逐段请求视频内容。

如果整个处理过程执行成功，资源服务会将分 P 文件的转码结果设置为成功，并写入文件路径、文件大小和播放时长等信息；如果处理过程中发生异常，例如临时文件不存在、合并失败、编码转换失败或切片失败，则将转码结果设置为失败。无论成功还是失败，资源服务都会在最终阶段调用主站服务接口回写处理结果，使主站服务能够根据所有分 P 文件的状态判断视频是否可以进入待审核阶段。

该实现方式能够较好地处理大文件资源。分片上传解决了单次上传压力问题，异步转码解决了投稿请求阻塞问题，FFmpeg 切片解决了大视频直接播放加载慢的问题，状态回写机制则保证了资源处理结果能够反馈到投稿业务流程中。

\subsection{HLS 播放资源生成与访问实现}

HLS 播放资源生成是资源服务面向视频播放功能提供支持的关键步骤。与直接返回完整 MP4 文件不同，HLS 会将视频拆分为一个索引文件和多个较小的切片文件。索引文件通常为 \texttt{m3u8} 格式，用于记录切片文件的播放顺序和时长；切片文件通常为 \texttt{ts} 格式，用于承载实际视频内容。用户播放视频时，播放器先请求 \texttt{m3u8} 索引文件，再按照索引内容逐段请求 \texttt{ts} 切片文件。

在资源生成阶段，系统通过 FFmpeg 命令完成 HLS 资源生成。资源服务先将完整视频转换为 \texttt{ts} 格式的中间文件，再使用分段命令按照固定时间间隔切分视频，并生成 \texttt{index.m3u8} 索引文件。切片时间设置为固定长度，可以使播放端按片段逐步加载视频，避免一次性加载完整视频文件造成较大带宽和内存压力。

HLS 播放资源生成完成后，系统会将播放资源所在目录保存到分 P 文件记录中。该目录通常与视频文件编号或上传路径相关联，目录下保存 \texttt{index.m3u8} 和多个 \texttt{ts} 切片。播放时，系统只需要根据分 P 文件编号请求资源服务，资源服务再根据文件编号解析出实际资源目录，并读取对应的索引文件或切片文件返回。

% 此处放图5-X：HLS 播放资源生成与访问流程图
% 建议图中包含：完整视频文件、FFmpeg 生成 ts、FFmpeg 切分 ts、生成 index.m3u8、请求 /videoResource/fileId、资源服务返回 m3u8、继续请求 ts、资源服务返回 ts 切片、记录播放信息。
% 图名建议：HLS 播放资源生成与访问流程图

播放资源访问主要包括索引文件访问和切片文件访问两个过程。索引文件访问根据 \texttt{fileId} 查询分 P 文件路径，并读取该目录下的 \texttt{m3u8} 文件返回；切片文件访问在 \texttt{fileId} 基础上额外携带具体的 \texttt{ts} 文件名，资源服务根据同一资源目录读取对应切片文件并返回。通过这种路径设计，系统可以避免直接暴露服务器真实文件路径，播放访问者只需要关心分 P 文件编号和切片名称即可。

资源服务在解析文件编号时，会优先查询正式视频文件记录。如果正式发布表中存在该文件，说明视频已经进入公开视频播放阶段，资源服务返回正式播放目录，并在读取索引文件时记录播放信息；如果正式视频文件不存在，系统会继续查询投稿态视频文件记录。该处理可以兼容审核前的资源预览场景，使创作者或审核流程能够访问投稿态资源，同时又能区分正式播放和非正式预览。

在正式视频播放场景下，资源服务读取 \texttt{m3u8} 索引文件后，会构造播放信息对象，并将视频编号、分 P 序号和用户编号等信息写入 Redis 播放记录队列。对于未登录用户，播放记录中可以只包含视频编号和分 P 信息；对于已登录用户，还可以记录用户编号。后续统计任务可以根据这些播放记录异步更新视频播放量，避免每次播放请求都直接修改数据库。

为了保证资源访问安全，系统在读取普通资源时需要校验资源路径是否合法，防止非法路径访问服务器其他目录。资源读取过程通过统一文件读取方法完成，该方法根据系统配置的项目文件根目录拼接资源路径，并将文件内容写入响应输出流。如果目标文件不存在，资源服务不会继续读取，避免发生文件访问异常。

综上，HLS 播放资源生成与访问实现完成了从完整视频文件到流式播放资源的转换。资源生成阶段通过 FFmpeg 输出 \texttt{m3u8} 索引和 \texttt{ts} 切片，资源访问阶段通过文件编号定位资源目录并返回播放文件，播放统计阶段通过 Redis 记录播放行为。该设计既提升了大视频播放体验，也保证了资源服务与视频业务状态之间能够保持一致。


\section{互动功能模块实现}

互动功能模块用于支撑用户在视频观看过程中的社区行为，主要包括评论、弹幕、点赞、收藏、投币和关注等功能。该模块与视频播放业务关系紧密，但其数据并不直接写入视频主数据表，而是通过独立的互动服务进行处理。互动服务负责保存用户行为记录、评论数据和弹幕数据，并根据行为结果更新视频相关统计字段。

从业务边界看，互动功能模块主要处理两类数据：一类是内容型互动数据，例如评论和弹幕；另一类是行为型互动数据，例如点赞、收藏、投币和关注。内容型互动数据需要保存用户输入的文本内容、所属视频、所属分 P 或父评论等信息；行为型互动数据则需要保存用户编号、视频编号、行为类型和行为次数等信息。系统通过这种拆分方式，使评论弹幕和行为记录能够分别维护，同时又能共同参与视频热度和用户活跃度统计。

互动功能模块中的大部分写操作都要求用户处于登录状态。系统在发表评论、发送弹幕、点赞、收藏、投币和关注之前，需要先从登录态中解析用户编号，并校验当前用户身份是否有效。未登录用户可以读取公开视频的评论和弹幕数据，但不能提交新的互动内容。通过这种方式，系统既支持游客浏览公开视频内容，又能保证用户行为与具体账号绑定。

% 此处放图5-X：互动功能模块协作图
% 建议图中包含：用户请求、网关服务、互动服务、评论服务、弹幕服务、用户行为服务、主站视频服务、MySQL、Redis。
% 图名建议：互动功能模块协作图

互动功能模块与主站视频服务之间存在跨服务协作关系。评论和弹幕在读取时，需要根据视频的互动设置判断该视频是否关闭评论区或弹幕功能；用户行为在写入后，需要影响视频的点赞数、收藏数、投币数、评论数或弹幕数等统计字段。因此，互动服务虽然独立处理互动数据，但仍需要与视频主数据保持一致。

\subsection{评论与弹幕实现}

评论功能用于支持用户围绕视频整体内容进行讨论。系统在发表评论时，需要接收视频编号、评论内容、回复评论编号和可选图片路径等参数。服务端首先校验用户登录状态，然后校验评论内容长度和视频编号有效性。校验通过后，系统构造评论对象，将用户编号、用户昵称、头像、视频编号、评论内容和回复关系写入评论数据表。

评论读取主要根据视频编号进行查询。系统会先获取视频对应的互动设置，如果该视频关闭评论功能，则直接返回空评论结果；如果评论功能开启，则按照评论排序规则查询评论数据。顶级评论通常按照点赞数或发布时间排序，回复评论则通过父评论编号与顶级评论形成层级关系。为了支持评论区置顶功能，系统还需要查询当前视频是否存在置顶评论，并在返回评论列表时将置顶评论放在较前位置。

对于已登录用户，评论读取结果还需要包含当前用户对评论的点赞或点踩状态。系统会根据当前用户编号、视频编号和评论行为类型查询用户行为记录，并将用户评论行为状态与评论列表一起返回。这样可以避免用户重复执行相同评论行为，也能在后续评论点赞或点踩操作中保持状态一致。

弹幕功能与评论功能不同，弹幕与具体视频分 P 文件和播放时间点绑定。用户发送弹幕时，需要提交视频编号、文件编号、弹幕内容、弹幕颜色、显示模式和播放时间。系统在保存弹幕前，首先校验用户登录状态，然后校验弹幕内容、视频编号、文件编号和播放时间是否有效。校验通过后，系统将弹幕数据写入弹幕表，并记录弹幕发送时间。

弹幕读取时，系统根据视频编号和文件编号查询对应弹幕数据，并按照弹幕编号或播放时间排序返回。由于一个视频可能包含多个分 P 文件，因此弹幕不能只按照视频编号查询，否则不同分 P 的弹幕会混合显示。通过视频编号和文件编号共同限定弹幕范围，可以保证弹幕内容与当前播放文件保持一致。

% 此处放图5-X：评论与弹幕处理流程图
% 建议图中包含：读取视频互动设置、判断评论/弹幕是否关闭、加载评论或弹幕、登录用户读取互动状态、发表评论、发送弹幕、写入评论表或弹幕表。
% 图名建议：评论与弹幕处理流程图

评论和弹幕实现的关键区别在于业务绑定对象不同。评论绑定视频整体内容，主要用于长文本讨论和回复；弹幕绑定具体分 P 文件和播放时间点，主要用于观看过程中的实时内容表达。因此，评论数据需要维护父子评论关系和置顶状态，而弹幕数据需要维护文件编号、播放时间、颜色和显示模式等播放相关字段。

\subsection{点赞、收藏、投币与关注实现}

点赞、收藏、投币与关注属于行为型互动功能。这类功能不以用户输入文本为主要内容，而是通过行为类型记录用户与视频或用户之间的关系。系统在处理这些操作时，首先需要校验用户登录状态，然后根据用户编号、视频编号、行为类型和可选评论编号构造用户行为对象，最后交由用户行为服务完成保存。

点赞功能主要用于表达用户对视频或评论的认可。视频点赞以视频编号和用户编号作为唯一行为约束，评论点赞还需要额外绑定评论编号。系统在保存点赞行为时，需要判断该用户是否已经执行过相同点赞操作。如果已经存在相同行为，则可以根据业务规则取消行为或拒绝重复写入；如果不存在，则新增用户行为记录，并更新对应视频或评论的点赞数量。

收藏功能用于记录用户希望后续再次访问的视频。收藏行为与用户编号和视频编号绑定，保存后可以在用户中心或个人收藏列表中查询。收藏功能的实现重点在于防止重复收藏，并在收藏或取消收藏时同步维护视频收藏数。由于收藏属于用户个人资产类数据，必须要求用户登录后才能执行。

投币功能用于表达用户对视频创作者的支持。与点赞和收藏相比，投币操作通常还涉及用户资源扣减和视频投币统计增加。因此，系统在处理投币前，需要校验投币数量是否合法，并判断用户是否具有足够资源。投币成功后，系统写入投币行为记录，并更新视频投币数量。为了避免重复扣减或统计异常，投币操作需要保证行为记录和统计更新在业务上保持一致。

关注功能用于建立用户与创作者之间的关系。当用户关注视频作者时，系统需要保存关注用户和被关注用户之间的关系；当用户取消关注时，系统需要删除或更新该关注关系。关注行为不直接绑定某一个视频，但通常由视频播放页或用户信息入口触发。系统在关注前需要判断关注对象是否合法，并避免用户关注自己。

% 此处放图5-X：点赞、收藏、投币与关注处理流程图
% 建议图中包含：用户发起互动操作、校验登录态、判断行为类型、校验是否重复、保存或取消行为记录、更新视频/评论/用户统计、返回操作结果。
% 图名建议：用户行为处理流程图

这些行为虽然表现形式不同，但底层实现具有相似性，均可以抽象为“用户编号、业务对象编号、行为类型、行为数量和行为时间”的记录结构。通过统一用户行为表保存点赞、收藏、投币和评论点赞等行为，可以减少不同互动功能之间的重复设计，同时便于后续统计用户活跃度和视频热度。

\subsection{互动统计数据更新实现}

互动统计数据更新用于保证视频详情页、首页列表和用户中心中展示的播放量、点赞数、收藏数、投币数、评论数和弹幕数等数据能够反映最新业务状态。由于互动行为频繁发生，如果每次行为都进行复杂查询，可能会增加数据库压力。因此，系统需要根据不同互动类型选择合适的统计更新方式。

评论统计主要在评论发布或删除时更新。用户发表评论成功后，系统可以增加视频评论数量；当用户删除评论时，系统需要减少对应评论数量，并处理该评论下的回复数据。对于评论点赞和点踩，则需要更新评论自身的点赞数或点踩数，而不是直接影响视频主表中的统计字段。

弹幕统计主要在弹幕发送成功后更新。弹幕与具体视频文件绑定，但在视频展示中通常以视频维度展示弹幕总数。因此，系统保存弹幕后，可以根据视频编号更新视频弹幕数量，也可以通过异步统计方式汇总所有分 P 文件下的弹幕数量。对于多分 P 视频，统计逻辑需要保证不同文件弹幕能够汇总到同一视频维度。

点赞、收藏和投币统计主要与用户行为记录关联。系统在保存行为记录前，需要判断该行为是否已经存在；只有当新增行为有效时，才更新对应统计数量。对于取消点赞或取消收藏场景，则需要同步减少统计数量。投币操作通常不允许简单撤销，因此统计更新更偏向单向增加，并与用户资源扣减保持一致。

关注统计主要作用于用户关系。用户关注创作者后，需要增加被关注用户的粉丝数，并增加当前用户的关注数；取消关注时则进行相反更新。该类统计虽然不直接属于视频主数据，但会影响作者主页、视频作者信息和用户中心展示，因此也属于互动统计数据的一部分。

% 此处放图5-X：互动统计数据更新流程图
% 建议图中包含：互动行为写入、判断统计对象、更新视频统计、更新评论统计、更新用户关注统计、返回最新状态。
% 图名建议：互动统计数据更新流程图

互动统计数据更新的关键是保证“行为记录”和“统计字段”之间的一致性。行为记录用于描述用户是否执行过某项操作，统计字段用于快速展示聚合结果。如果只保存行为记录而不维护统计字段，则列表和详情页展示需要频繁聚合查询；如果只更新统计字段而不保存行为记录，则无法判断用户是否已点赞、收藏或关注。因此，系统需要同时维护行为明细和统计结果。

整体来看，互动功能模块通过评论表、弹幕表、用户行为表和统计字段共同完成社区互动能力。评论和弹幕提供内容表达，点赞、收藏、投币和关注提供行为反馈，统计数据更新则将用户行为转化为视频热度和用户关系数据。该模块为视频分享平台提供了基本的社区交互能力。

\section{后台管理功能实现}

后台管理功能面向系统管理员提供分类维护、内容管理和视频审核等能力。与普通用户端功能不同，后台管理功能更强调数据治理、内容安全和系统运营。管理员可以维护视频分类结构，查看投稿视频列表，读取视频投稿详情，查看 AI 审核结果，并对待审核视频执行人工通过或驳回操作。

后台管理服务通过独立的管理员认证机制保护管理接口。管理员登录时需要校验验证码、后台账号和后台密码，登录成功后系统生成后台管理令牌并写入 Redis。后续访问后台接口时，系统通过后台令牌判断管理员身份。通过区分普通用户登录态和管理员登录态，可以避免普通用户访问后台管理能力。

从服务协作角度看，后台管理端自身主要提供管理入口，部分业务处理需要调用主站业务服务完成。例如，视频审核列表、审核详情、AI 审核摘要、AI 审核明细和人工审核操作由后台管理服务接收请求后，通过服务调用转发到主站业务服务处理；分类管理则由后台管理服务调用分类业务服务直接维护分类数据。该设计使后台管理模块能够复用主站已有的视频和分类业务能力。

% 此处放图5-X：后台管理功能整体协作图
% 建议图中包含：管理员、后台管理服务、管理员认证、分类管理、视频内容管理、视频审核管理、主站业务服务、AI审核数据、MySQL、Redis。
% 图名建议：后台管理功能整体协作图

后台管理功能的实现重点在于权限控制和状态控制。权限控制用于保证只有管理员可以访问管理接口；状态控制用于保证视频审核、分类删除和内容管理等操作符合业务规则。例如，视频只有在转码完成并进入待审核状态后才适合进行人工审核，分类删除需要考虑是否存在子分类或视频引用，审核通过后需要将投稿态数据迁移到正式发布数据。

\subsection{后台管理端整体实现}

后台管理端整体实现主要包括管理员认证、管理接口转发、业务数据查询和管理操作执行几个部分。管理员在访问后台管理功能前，需要先通过验证码和管理员账号密码校验。系统使用验证码防止恶意重复登录请求，使用 Redis 保存管理员登录令牌，使用 Cookie 或请求标识维持后续管理会话。

管理员登录成功后，可以访问后台分类管理、视频审核管理和内容管理等接口。系统在处理这些接口之前，需要判断当前请求是否携带有效的后台登录令牌。若后台登录态无效，则拒绝访问管理功能；若登录态有效，则进入对应管理业务处理流程。通过统一后台认证，可以避免每个管理接口重复编写管理员身份校验逻辑。

后台管理端的业务处理并不完全依赖自身服务完成。对于视频投稿、AI 审核结果和人工审核操作，后台管理服务通过内部服务调用访问主站业务服务。主站业务服务负责查询投稿视频、审核状态、AI 审核任务和分 P 审核明细，并在管理员执行审核操作时更新投稿状态。后台管理服务在这个过程中承担管理入口和权限控制作用。

% 此处放图5-X：后台管理端请求处理流程图
% 建议图中包含：管理员请求、校验后台登录态、判断管理功能类型、分类管理本地处理、视频管理调用主站服务、返回管理结果。
% 图名建议：后台管理端整体处理流程图

后台管理端整体实现的关键在于将管理员操作与普通用户操作隔离。普通用户只能访问公开视频、投稿和互动功能，而管理员能够访问全站分类、投稿视频审核和 AI 审核结果等系统级数据。因此，后台管理接口需要使用独立路径、独立登录令牌和独立访问控制逻辑，避免用户端权限和管理端权限混用。

\subsection{分类与内容管理实现}

分类管理用于维护视频内容的分类体系。系统中的视频分类通常采用一级分类和二级分类结构，一级分类用于划分较大的内容领域，二级分类用于进一步细化视频类型。后台管理模块提供分类读取、分类保存、分类删除和分类排序等功能，管理员可以根据平台运营需求调整分类结构。

分类读取时，系统按照排序字段升序查询分类数据，并将查询结果转换为树形结构返回。树形结构能够表达父分类和子分类之间的层级关系，使分类管理和用户端分类浏览使用统一的数据来源。分类保存时，系统需要接收父分类编号、分类编号、分类编码、分类名称和图标等信息；如果分类编号为空，则表示新增分类；如果分类编号存在，则表示修改已有分类。

分类删除需要进行业务校验。若待删除分类下存在子分类，或者已有视频引用该分类，则直接删除可能导致视频分类数据异常。因此，系统在删除分类前应校验分类是否仍被使用。只有满足删除条件时，才允许删除分类记录。分类排序则根据同一父分类下提交的分类编号顺序更新排序字段，使分类展示顺序能够按照管理员配置生效。

内容管理主要围绕视频投稿数据展开。后台管理端可以按照投稿状态、视频名称、用户信息或更新时间等条件查询投稿视频列表。列表数据不仅包含投稿标题、封面、分类、投稿用户和状态，还可以包含播放、点赞、弹幕、评论、投币和收藏等统计数据。对于已经接入 AI 审核结果的投稿，列表还可以展示 AI 建议、风险等级和摘要信息，用于辅助管理员快速判断内容风险。

% 此处放图5-X：分类与内容管理流程图
% 建议图中包含：管理员进入管理功能、加载分类树、保存分类、删除分类、调整排序、查询投稿内容列表、查看内容详情。
% 图名建议：分类与内容管理流程图

分类与内容管理的实现目标是为平台提供基础运营能力。分类管理保证视频内容能够按照合理结构组织，内容管理保证管理员能够从全站角度查看投稿状态和内容质量。二者共同为视频审核和内容发布提供基础数据支持。

\subsection{视频审核管理实现}

视频审核管理是后台管理功能中的核心业务。视频从创作者投稿到正式发布，需要经过上传、转码、待审核、人工审核和正式发布等阶段。资源服务完成全部分 P 文件转码后，主站服务会将投稿状态更新为待审核，并生成或关联 AI 审核任务。管理员在后台管理中查看待审核视频列表，根据视频内容、投稿信息和 AI 审核结果进行最终审核。

后台审核列表通过后台管理服务调用主站业务服务获取投稿视频数据。系统根据查询条件读取投稿表中的视频记录，并关联用户信息、视频统计信息和 AI 审核摘要。管理员可以通过列表快速查看视频标题、封面、投稿用户、投稿状态、更新时间以及 AI 审核建议等内容。如果需要进一步查看详情，系统可以根据视频编号读取投稿详情、分 P 文件信息、AI 视频级审核摘要和 AI 分 P 审核明细。

AI 审核结果在本系统中主要起辅助决策作用。AI 审核任务表用于保存视频级审核结果，例如 AI 建议、风险等级、审核摘要、模型名称、模型版本、任务状态和完成时间等；AI 审核明细表用于保存分 P 文件级审核结果，例如分 P 文件编号、风险分值、风险标签、命中片段和审核说明等。管理员可以结合这些数据判断视频是否存在违规风险。

% 此处放图5-X：视频审核管理流程图
% 建议图中包含：加载待审核视频列表、查看投稿详情、读取AI审核摘要、读取AI分P明细、管理员判断、审核通过、审核驳回、更新投稿状态、通过时迁移到正式视频表。
% 图名建议：视频审核管理流程图

人工审核操作分为通过和驳回。管理员选择通过时，系统将投稿状态更新为审核通过，并将投稿态视频主数据和分 P 文件数据同步到正式视频表和正式视频文件表中。正式表中的视频随后可以进入首页列表、搜索索引和公开视频播放流程。审核通过后，系统还需要处理旧文件清理、文件更新标识重置和搜索索引写入等后续操作，保证正式数据和播放资源保持一致。

管理员选择驳回时，系统将投稿状态更新为审核不通过，并保留投稿态数据供创作者后续查看或重新编辑。驳回的视频不会迁移到正式视频表，也不会进入公开视频列表和搜索结果。若系统支持驳回原因，则应将驳回理由与投稿记录关联保存，使创作者能够根据原因修改后重新提交。

视频审核状态流转需要严格控制。转码中的视频不能直接审核，因为播放资源尚未生成；转码失败的视频需要创作者重新上传或编辑；待审核视频可以进入人工审核；审核通过后进入正式发布状态；审核不通过后回到创作者修改流程。通过明确状态边界，系统能够避免未完成处理的视频被发布，也能避免审核结果与资源状态不一致。

% 此处放图5-X：视频审核状态流转图
% 建议图中包含：转码中、转码失败、待审核、AI审核完成、人工审核通过、人工审核驳回、正式发布、重新编辑。
% 图名建议：视频审核状态流转图

综上，视频审核管理模块通过投稿状态、AI 审核结果和人工审核操作共同完成内容发布控制。AI 审核负责提供风险提示和辅助判断，管理员负责最终审核决策，主站业务服务负责根据审核结果迁移或保留投稿数据。该设计能够在保证内容安全的同时，使视频发布流程具有明确的数据状态和操作边界。

\section{AI 审核任务编排与人机协同实现}

AI 审核任务编排位于视频转码完成之后，是连接资源处理、AI 初审和人工终审的实现环节。本节说明主站服务在代码层面如何创建 AI 审核任务、投递审核消息、接收审核结果，并将 AI 初审结果提供给管理员进行最终判断。

在实现上，AI 审核并不直接决定视频是否发布。系统将 AI 审核定位为自动化初审环节，负责识别风险片段、生成风险等级和审核建议；人工审核则作为最终审核环节，负责结合平台规则和视频内容作出通过或驳回决定。通过这种方式，系统既可以利用 AI 降低人工筛查成本，又可以避免模型误判直接影响视频发布结果。

% 图5-x：AI 审核任务编排活动图
% 建议图类型：活动图
% 建议内容：
% 转码完成 -> 创建 AI 审核任务 -> 创建分 P 审核明细 -> 投递审核请求消息
% -> AI 服务处理视频 -> 返回审核结果 -> 更新 AI 审核任务状态
% -> 后台展示 AI 结果 -> 管理员人工终审。
% 该图用于说明 AI 审核任务在主站服务、消息队列、AI 服务和后台审核之间的流转关系。


\subsection{AI 审核任务创建与消息投递实现}

当资源服务完成视频分 P 文件转码后，主站服务会根据投稿编号查询该投稿下所有分 P 文件的处理状态。只有所有分 P 文件均已生成可审核的视频资源路径时，系统才会创建 AI 审核任务。这样可以避免 AI 服务收到不完整的视频文件信息，从而导致审核失败或审核结果缺失。

AI 审核任务创建时，系统首先生成唯一的审核请求编号，用于标识本次审核流程。随后，主站服务在 AI 审核任务主表中写入视频级任务记录，并根据投稿下的分 P 文件列表创建分 P 审核明细记录。视频级任务用于描述本次审核的整体状态，分 P 明细用于记录每个视频文件的审核结果、风险标签和命中片段。

消息投递时，系统不会将视频文件本身写入消息队列，而是将视频编号、审核任务编号、分 P 文件编号、文件路径和分 P 序号等信息封装为审核请求消息。AI 审核服务消费该消息后，根据文件路径读取待审核视频并执行多模态审核。该方式能够避免大文件通过消息队列传输，同时也便于通过数据库记录追踪审核任务状态。

% 图5-x：AI 审核任务创建与消息投递时序图
% 建议图类型：时序图
% 建议参与对象：
% 资源服务、主站服务、数据库、RabbitMQ、AI 审核服务。
% 建议流程：
% 1. 资源服务通知转码完成；
% 2. 主站服务查询投稿和分 P 文件；
% 3. 主站服务写入 ai_audit_task；
% 4. 主站服务批量写入 ai_audit_task_item；
% 5. 主站服务向 RabbitMQ 投递 audit.video.request；
% 6. AI 审核服务消费审核请求。
% 该图用于替代表格说明任务创建和消息投递过程。

该实现的重点是保证“任务先落库，消息后投递”。系统先写入审核任务数据，再发送审核请求消息。若消息投递成功，则任务状态更新为处理中；若消息投递失败，则记录失败原因，并将任务状态更新为失败或等待重试。这样即使 AI 服务暂时不可用，系统也能够通过任务表定位异常任务，不会造成审核请求完全丢失。


\subsection{审核结果接收与任务状态更新实现}

AI 审核服务完成视频分析后，会将审核结果发送回主站服务。主站服务消费审核结果消息时，首先根据审核请求编号查询对应的 AI 审核任务，并校验视频编号和审核版本。若当前结果消息对应的是过期任务，系统不应覆盖新的审核结果，以避免编辑重审场景下旧结果污染当前数据。

审核结果分为视频级结果和分 P 级结果。视频级结果用于更新 AI 审核任务主记录，包括 AI 建议、风险等级、审核摘要和完成时间等信息；分 P 级结果用于更新 AI 审核任务明细，包括分 P 审核建议、风险分值、风险标签、命中片段和审核原因等信息。命中片段通常保存为结构化 JSON，记录风险出现的开始时间、结束时间、代表帧路径、文本摘要、声音事件和模型判断理由。

% 图5-x：AI 审核任务状态流转图
% 建议图类型：状态图
% 建议状态：
% 待处理、处理中、已完成、失败、取消。
% 建议转移关系：
% 创建任务 -> 待处理；
% 消息投递成功 -> 处理中；
% AI 结果回写成功 -> 已完成；
% 消息投递失败或结果解析失败 -> 失败；
% 被新审核版本替代或无有效文件 -> 取消。
% 该图用于替代表格说明 AI 审核任务状态变化。

在结果回写过程中，系统需要保证主任务和分 P 明细之间的状态一致性。只有当所有分 P 审核明细均完成更新后，视频级 AI 审核任务才可以标记为已完成。如果部分分 P 结果解析失败或写入失败，则任务应标记为失败，并保留异常原因，便于后续人工排查或重新触发审核。

通过该实现，AI 服务输出的模型结果被转换为可查询、可展示、可复核的结构化审核数据。后台审核页面不需要再次调用 AI 服务，而是直接读取数据库中的审核任务和分 P 明细，从而提升后台审核页面的响应效率。


\subsection{AI 初审与人工终审协同实现}

AI 初审与人工终审的协同关系主要体现在审核权限边界上。AI 服务负责对视频内容进行自动化分析，输出风险等级、审核建议和命中证据；管理员负责查看 AI 初审结果，并结合视频内容和平台审核规则作出最终判断。AI 不直接修改正式发布数据，也不直接决定视频通过或驳回。

在后台审核详情页中，系统需要同时展示投稿基础信息、视频播放内容和 AI 审核结果。投稿基础信息包括标题、分类、封面、标签和简介；视频播放内容用于管理员直接查看原始视频；AI 审核结果包括 AI 建议、风险等级、风险摘要、命中片段和代表帧。管理员可以优先查看 AI 命中的风险片段，再决定是否完整播放视频进行复核。

% 图5-x：AI 初审与人工终审协同流程图
% 建议图类型：流程图
% 建议内容：
% AI 初审完成 -> 后台展示 AI 结果 -> 管理员查看命中片段
% -> 管理员人工判断 -> 审核通过 / 审核驳回。
% 图中应突出 AI 只提供建议，人工审核才是最终决策节点。

为了降低人工审核成本，系统将 AI 命中的片段作为重点证据保存。每个命中片段包含风险类型、时间范围、文本摘要、声音事件、代表帧路径和判断原因。管理员可以根据命中时间点快速定位问题片段，而不需要从头完整观看所有视频内容。

该设计使 AI 审核和人工审核形成互补关系。AI 审核提高风险发现效率，人工审核保证最终结论可靠性。对于低风险内容，管理员可以快速复核；对于中高风险内容，管理员可以根据命中片段重点判断；对于模型证据不足的内容，系统可以将其标记为需要人工复核。


\subsection{审核通过与发布数据同步实现}

当管理员执行审核通过操作后，系统需要将投稿态数据同步为正式发布数据。投稿态数据主要用于支撑上传、转码和审核流程，正式发布数据则用于支撑前台视频浏览、搜索和播放。系统通过投稿态和发布态分离，保证未审核通过的视频不会进入正式播放链路。

审核通过时，系统首先更新投稿主记录状态，将其标记为审核通过。随后，系统将投稿主信息写入或更新到正式视频表，包括标题、封面、分类、标签、简介、作者和总时长等字段。对于分 P 文件，系统将投稿分 P 表中的文件名称、文件路径、播放资源路径、文件序号和时长等信息同步到正式分 P 文件表中，使前台播放页能够读取正式视频资源。

% 图5-x：审核通过与发布数据同步流程图
% 建议图类型：流程图
% 建议内容：
% 管理员点击通过 -> 更新 video_info_post 状态
% -> 同步投稿主信息到 video_info
% -> 同步分 P 文件到 video_info_file
% -> 重置分 P 更新标识
% -> 更新搜索索引或播放数据。
% 如果管理员驳回，则只更新投稿状态和驳回原因，不写入正式发布表。

当管理员执行驳回操作时，系统只更新投稿状态和驳回原因，不会将数据同步到正式视频表。这样可以保证审核不通过的视频不会出现在首页、搜索、推荐和播放页面中。通过这种方式，系统将 AI 审核、人工审核和发布同步三个环节解耦，既能保留完整审核记录，也能保证前台内容发布的安全性。


\section{多模态 AI 审核服务实现}

多模态 AI 审核服务是 AI 初审的具体执行模块。该服务接收主站服务投递的视频审核请求后，按照分 P 文件逐个处理视频内容，并输出分 P 级和视频级审核结果。与第二章中对 Whisper、YAMNet 和 Qwen-VL 等技术的介绍不同，本节重点说明这些模型在系统中的调用顺序、数据转换方式和结果汇总方式。

AI 审核服务不会直接将完整视频输入多模态模型，而是先对视频进行预处理，将其拆分为音频、语音文本、声音事件和关键帧等结构化证据。随后，系统以时间窗口为单位构建审核片段，将同一时间段内的画面、语音文本、声音事件和投稿元数据组合为 \texttt{segment}，再调用多模态模型进行审核。


\subsection{AI 审核服务结构}

AI 审核服务可以划分为审核编排层、视频处理层、片段构建层和结果生成层。审核编排层负责控制完整审核流程；视频处理层负责完成音频提取、语音识别、声音事件识别和关键帧抽取；片段构建层负责将不同模态的结果按照时间戳对齐；结果生成层负责将多个片段结果汇总为最终审核结论。

% 图5-x：AI 审核服务核心类图
% 建议图类型：类图
% 建议类：
% ModerationEngine
% VideoModerationService
% AuditSegment
% AuditResult
% 建议关系：
% ModerationEngine 依赖 VideoModerationService；
% VideoModerationService 负责 extract_audio、detect_voice_segments、
% detect_sound_events、transcribe_audio_segments、extract_keyframes、
% build_audit_segments、batch_audit 等方法；
% AuditSegment 聚合语音文本、声音事件、关键帧和审核结果；
% AuditResult 保存风险类型、风险等级、审核建议和原因说明。
% 该图用于说明 AI 审核服务内部对象之间的职责关系。

其中，\texttt{ModerationEngine} 是审核流程的入口对象，负责任务级流程编排；\texttt{VideoModerationService} 封装具体的视频处理和模型调用能力；\texttt{segment} 是审核过程中的核心数据结构，用于表示视频时间轴上的一个审核窗口；\texttt{audit\_result} 用于保存片段级模型审核结果。

这种结构将流程控制和模型能力调用分离。审核编排层只负责控制处理顺序和结果流转，具体的音频处理、关键帧抽取、声音事件识别和多模态模型调用封装在服务类中，从而降低代码耦合度，也便于后续替换或扩展模型能力。


\subsection{视频预处理与音频提取实现}

视频预处理的目标是将原始视频文件转换为后续审核模型可以使用的数据。系统首先检测视频是否存在音频流。如果视频存在音频流，则使用 FFmpeg 将音频提取为 16kHz 的 WAV 文件，供 VAD、Whisper 和 YAMNet 使用；如果视频不存在音频流，则跳过音频分析，但仍继续执行关键帧抽取和视觉审核，避免纯画面视频被漏审。

音频提取完成后，系统会形成两条分析路径。第一条是语音分析路径，即通过 VAD 检测有效语音片段，再调用 Whisper 将语音内容转换为文本。第二条是声音事件分析路径，即通过 YAMNet 识别音频中的环境声音或异常声音。二者的作用不同：Whisper 用于获得“说了什么”，YAMNet 用于判断“出现了什么声音”。

为了降低 VAD 漏检带来的风险，系统不应将 VAD 作为 Whisper 的唯一入口。当 VAD 检测结果为空，或语音覆盖率明显偏低时，可以启动整段音频转写兜底机制。这样即使 VAD 未能正确检测出语音段，音频中的违规语义仍有机会被转写并进入后续审核。

% 图5-x：视频预处理活动图
% 建议图类型：活动图
% 建议内容：
% 输入视频 -> 检测音频流
% -> 有音频：提取 wav -> VAD -> Whisper -> YAMNet
% -> 无音频：跳过音频分析
% -> 提取关键帧 -> 输出语音文本、声音事件和关键帧。
% 由于本小节已经包含算法，该图可以选择性保留；如果第五章图较多，可以不画。

视频预处理与音频提取过程如算法~\ref{alg:video-preprocess-impl} 所示。

\begin{algorithm}[H]
    \caption{视频预处理与音频提取实现算法}
    \label{alg:video-preprocess-impl}
    \begin{algorithmic}[1]
        \Require 视频文件路径 $videoPath$
        \Ensure 语音文本集合 $S$，声音事件集合 $A$，关键帧集合 $F$
        \State 检测视频文件是否存在音频流
        \If{存在音频流}
            \State 使用 FFmpeg 提取 16kHz WAV 音频文件
            \State 使用 VAD 检测语音片段集合 $V$
            \State 使用 YAMNet 检测声音事件集合 $A$
            \If{$V$ 不为空}
                \State 调用 Whisper 对语音片段进行分段转写，得到语音文本集合 $S$
            \Else
                \State 调用 Whisper 对整段音频进行兜底转写，得到语音文本集合 $S$
            \EndIf
        \Else
            \State 设置 $S \leftarrow \emptyset$，$A \leftarrow \emptyset$
        \EndIf
        \State 根据固定间隔、画面变化和风险声音时间点提取关键帧集合 $F$
        \State 返回 $S$、$A$ 和 $F$
    \end{algorithmic}
\end{algorithm}

通过上述处理，系统将原始视频转换为语音文本、声音事件和关键帧三类证据。后续审核不再直接依赖完整视频文件，而是围绕这些带时间戳的结构化结果构建审核片段。


\subsection{语音文本、声音事件与关键帧审核实现}

完成视频预处理后，系统需要将语音文本、声音事件和关键帧组织为审核片段。审核片段是 AI 审核服务中的核心输入单位。每个片段对应视频时间轴上的一个时间窗口，系统将该窗口内的语音文本、声音事件和代表帧组合起来，形成可送入 Qwen 审核的多模态输入。

在文本处理方面，Whisper 的转写结果会被保存为带时间戳的文本片段。系统根据文本片段的开始时间和结束时间创建文本审核窗口，并在窗口内查找对应的关键帧和声音事件。对于没有语音文本的视频，系统不能直接跳过，而是需要基于关键帧创建视觉审核窗口。

在声音事件处理方面，YAMNet 的识别结果会作为片段上下文加入审核输入。对于普通声音事件，系统仅将其作为辅助说明；对于枪声、爆炸、尖叫、警报等风险声音，系统会以声音事件时间点为中心创建声音事件审核窗口。如果该时间点附近能够找到关键帧，则将关键帧一起送入多模态模型；如果无法找到可用画面，则将该片段标记为需要人工复核。

在视觉处理方面，系统不将所有视频帧直接送入模型，而是提取关键帧。关键帧来源包括固定间隔采样帧、画面变化帧和风险声音附近补充帧。这样可以在控制模型调用次数的同时，提高短时风险画面的覆盖率。

% 图5-x：多模态审核片段构建流程图
% 建议图类型：流程图
% 建议内容：
% 语音文本片段 -> 构建 text-visual segment；
% 未覆盖关键帧 -> 构建 visual-only segment；
% 风险声音事件 -> 构建 audio-event segment；
% segment 集合 -> Qwen 多模态审核。
% 该图用于突出 segment 不再只依赖语音，而是由文本、声音事件和关键帧共同构成。

语音文本、声音事件与关键帧的片段构建过程如算法~\ref{alg:segment-build-impl} 所示。

\begin{algorithm}[H]
    \caption{语音文本、声音事件与关键帧审核片段构建算法}
    \label{alg:segment-build-impl}
    \begin{algorithmic}[1]
        \Require 语音文本集合 $S$，声音事件集合 $A$，关键帧集合 $F$
        \Ensure 审核片段集合 $G$
        \State 初始化审核片段集合 $G \leftarrow \emptyset$
        \For{每一个语音文本片段 $s \in S$}
            \State 根据 $s.start$ 和 $s.end$ 创建文本审核窗口
            \State 查找该窗口内的关键帧和声音事件
            \State 生成包含语音文本、声音事件和视觉帧的 \texttt{text-visual segment}
            \State 将该片段加入 $G$
        \EndFor
        \For{每一个未被已有片段覆盖的关键帧 $f \in F$}
            \State 以 $f.timestamp$ 为中心创建视觉审核窗口
            \State 生成 \texttt{visual-only segment}
            \State 将该片段加入 $G$
        \EndFor
        \For{每一个风险声音事件 $a \in A$}
            \If{$a$ 未被已有片段覆盖}
                \State 以 $a.timestamp$ 为中心创建声音事件审核窗口
                \State 查找距离 $a.timestamp$ 最近的关键帧作为视觉证据
                \State 生成 \texttt{audio-event segment}
                \State 将该片段加入 $G$
            \EndIf
        \EndFor
        \State 按片段开始时间对 $G$ 排序
        \State 返回审核片段集合 $G$
    \end{algorithmic}
\end{algorithm}

构建完成后，每个 \texttt{segment} 会被送入 Qwen 审核。Qwen 的输入主要包括代表帧和文本提示词。提示词中包含当前片段的语音转写内容、声音事件摘要、投稿标题、简介、标签以及审核规则说明。由于 Qwen 不直接接收原始音频，因此语音和声音事件需要先转换为文本化信息，再与关键帧一起参与模型判断。

通过该实现，系统可以覆盖四类审核场景：第一类是含语音的视频片段，系统通过语音文本和关键帧共同判断；第二类是无语音但存在画面风险的视频片段，系统通过 \texttt{visual-only segment} 保证其进入审核；第三类是存在风险声音的视频片段，系统通过 \texttt{audio-event segment} 进行重点检查；第四类是 VAD 漏检的视频片段，系统通过整段音频兜底转写降低漏审风险。


\subsection{综合风险结果生成实现}

片段级审核完成后，系统需要将多个 \texttt{segment} 的审核结果汇总为分 P 级结果，再进一步汇总为视频级结果。分 P 级结果用于说明某个视频文件是否存在风险，视频级结果用于后台审核列表展示和管理员最终审核参考。

综合风险生成采用高风险优先原则。如果任一片段被判断为高风险，则该分 P 的建议结果应倾向于驳回；如果没有明确高风险，但存在模型判断不确定、缺少视觉证据或中风险内容，则建议人工复核；只有所有片段均未发现明显风险时，才建议通过。

% 图5-x：综合风险结果生成流程图
% 建议图类型：流程图
% 建议内容：
% 片段审核结果集合 -> 提取风险片段 -> 归并风险等级
% -> 生成分 P 级审核建议 -> 汇总视频级审核建议
% -> 返回主站服务。
% 不建议再使用算法，避免本章算法过多。

在生成分 P 结果时，系统会遍历所有片段审核结果，提取命中风险的片段，并记录其时间范围、风险类型、风险原因、代表帧路径和文本摘要。若存在高风险片段，则分 P 审核建议为驳回；若存在中风险或审核证据不足的片段，则分 P 审核建议为人工复核；若所有片段均为低风险，则分 P 审核建议为通过。

在生成视频级结果时，系统继续对所有分 P 结果进行归并。若任一分 P 建议驳回，则视频级 AI 建议为驳回；若不存在驳回但存在人工复核，则视频级 AI 建议为人工复核；只有所有分 P 均建议通过时，视频级 AI 建议才为通过。该策略能够避免高风险片段被其他低风险片段稀释，使审核结果更加谨慎。

最终，AI 服务将视频级结果和分 P 级结果封装为审核结果消息返回主站服务。主站服务将结果写入 AI 审核任务表和分 P 明细表，供管理员在后台审核页面查看。通过这种从片段结果到分 P 结果、再到视频级结果的逐级汇总方式，系统既能给出整体审核建议，又能保留模型判断的具体证据，便于人工复核。